% Counter reference of the register-spill study (metric schema 10): every
% counter the detector writes, the gem5 statistics the reports read, and every
% formula the reports compute, each with how to read it.
% Same preamble and author block as spill_definition.tex (metric schema 7).
% Section 12 repeats the worked example of section 10 in Turkish.
% Build:  tectonic general_reports/definitions/counter_reference_schema10.tex
\documentclass[10pt,a4paper]{article}
\usepackage[a4paper,top=16mm,bottom=16mm,left=20mm,right=20mm]{geometry}
\usepackage{booktabs}
\usepackage{tabularx}
\usepackage{xltabular}
\usepackage{lmodern}
\usepackage[T1]{fontenc}
\usepackage{amsmath}
\usepackage{microtype}
\usepackage{titlesec}

\pagestyle{empty}
\setlength{\parindent}{0pt}
\setlength{\parskip}{5pt}
\linespread{1.06}\selectfont
\raggedbottom

\titleformat{\section}{\normalfont\large\bfseries}{\thesection.}{0.6em}{}
\titlespacing*{\section}{0pt}{12pt plus 2pt minus 2pt}{5pt}
\titleformat{\subsection}{\normalfont\normalsize\bfseries}{\thesubsection}{0.6em}{}
\titlespacing*{\subsection}{0pt}{8pt plus 2pt minus 2pt}{3pt}

% Author-block markup kept in the IEEEtran spelling the paper uses, so the same
% block can be moved between documents unchanged. This report stays
% single-column: the tables need the full text width to stay legible.
% The affiliation block contains its own \\ line breaks. article's \author
% puts each author in a tabular, where a brace group cannot span a row break,
% so the affiliation goes into a nested tabular instead.
\newcommand{\IEEEauthorblockN}[1]{{\normalsize #1}}
\newcommand{\IEEEauthorblockA}[1]{%
  \\[2pt]{\small\begin{tabular}[t]{@{}c@{}}#1\end{tabular}}}

\newcommand{\LongDate}{%
    \number\day\space
    \ifcase\month
    \or January%
    \or February%
    \or March%
    \or April%
    \or May%
    \or June%
    \or July%
    \or August%
    \or September%
    \or October%
    \or November%
    \or December%
    \fi
    \space \number\year
}

% X columns whose share of the width can be tuned; C: a counter-name column.
\newcolumntype{L}[1]{>{\raggedright\arraybackslash\hsize=#1\hsize}X}
\newcolumntype{R}{>{\ttfamily}l}
\newcolumntype{C}[1]{>{\raggedright\ttfamily\arraybackslash}p{#1}}
\newcommand{\ct}[1]{\texttt{#1}}
% A breakable underscore for long counter names.
\newcommand{\us}{\_\allowbreak}

\setlength{\tabcolsep}{5pt}
\renewcommand{\arraystretch}{1.28}

\newcommand{\tablehead}[2]{%
  \nobreak{\small\textbf{Table #1.} #2}\par\nobreak\vspace{2pt}}

\title{
    \Large Improving CPU Efficiency for Data-Intensive Applications \\
    \Large through Cache-Aware Spill Optimization \\[6pt]
    \large Counter Reference: What Is Counted and How to Read It \\[5pt]
    \normalsize gem5 TimingSimpleCPU $\cdot$ RISC-V (RV64GC)
    $\cdot$ SPEC CPU2017 $\cdot$ metric schema 10
}

\author{
    \IEEEauthorblockN{L\"utfullah Burak Kaya}
    \IEEEauthorblockA{
        Ozyegin University \\
        burak.kaya.50711@ozu.edu.tr
    }
    \and
    \IEEEauthorblockN{Assistant Prof.\ Dr.\ \.{I}smail Akt\"urk}
    \IEEEauthorblockA{
        Ozyegin University \\
        ismail.akturk@ozyegin.edu.tr
    }
}

\date{\small \LongDate}

\begin{document}

\maketitle
\thispagestyle{empty}
\vspace{-9mm}

\section{What is measured}

A spill is detected \emph{dynamically}, from the runtime effective addresses
of the committed loads and stores alone. No compiler information, no register
name and no base-register test takes part. A store to the process stack is a
\emph{spill store}; a later load that reads a byte it wrote is a \emph{spill
load}; a spill store read back at least once is a \emph{spill event}.

This is an operational definition: \emph{runtime-detected stack store--load
reuse}, or spill-like memory traffic. Register-allocator spills,
callee-saved register saves and restores (return address included),
memory-resident locals and stack-passed arguments all meet it and cannot be
told apart from addresses alone. The counts are therefore not counts of
compiler-generated spill instructions.

Counting is confined to the region of interest (ROI), marked by
\ct{m5\_work\_begin} and \ct{m5\_work\_end} in the benchmark sources. In a
SimPoint run the ROI is the simulated interval after the warmup. Stores before
the ROI, the warmup included, are not tracked, so a reload of one never pairs.

\tablehead{1}{Counter names. Ratios are not counters: the report scripts
compute them. gem5's own statistics keep their names.}

\begin{tabularx}{\textwidth}{@{}C{52mm} X@{}}
\toprule
\normalfont\textbf{Name} & \textbf{Scope} \\
\midrule
total\_\textit{name} & The whole simulation; for a SimPoint estimate, the
  whole program (from the qemu profile). \\
roi\_\textit{name} & Only between the ROI markers; SimPoint: the interval. \\
diag\_\textit{name} & A technical diagnostic: checks the simulation and the
  matching, not a result. \\
roi\_\textit{p}\_\{l1,l2\}\_\{hits,misses\} & The L1D or L2 outcome of the
  accesses of population \textit{p} $\in$ \{\ct{spill\_store},
  \ct{spill\_event}, \ct{spill\_load}\}. \\
\bottomrule
\end{tabularx}

\section{When a store or a load counts}

\tablehead{2}{The conditions, checked when the instruction commits. S: a
store, L: a load.}

\begin{tabularx}{\textwidth}{@{}l L{1.30} L{0.70}@{}}
\toprule
& \textbf{Condition} & \textbf{Why} \\
\midrule
S1 & An ordinary store: not a store-conditional (SC), not an atomic memory
     operation (AMO)
   & Atomics are out of scope \\
S2 & It commits inside the ROI
   & Measurement window \\
S3 & Its effective address lies in the process stack,
     $\ct{stackMin} \le \text{addr} < \ct{stackBase}$
   & Spill slots live on the stack \\
S4 & It completes without a fault
   & A faulted attempt is discarded; the re-executed store counts once \\
\midrule
L1 & An ordinary load: not a load-reserved (LR)
   & As S1 \\
L2 & It commits inside the ROI
   & As S2 \\
L3 & Its effective address lies in the process stack
   & As S3 \\
L4 & Its start byte was last written by a store meeting S1--S4. The newest
     writer owns each byte, so after a partial overwrite the newer store owns
     the overlapped bytes; the widths may differ
   & The value read was produced by that store \\
L5 & The load is later than that store
   & Written before read \\
L6 & That store still owns the byte when the load commits
   & Not overwritten in between \\
L7 & It completes without a fault
   & As S4 \\
\bottomrule
\end{tabularx}

A store meeting S1--S4 is counted in \ct{roi\_spill\_store}, read back or
not. A load meeting L1--L7 is counted in \ct{roi\_spill\_load}; the first such
load of a store also counts that store once in \ct{roi\_spill\_event}. A
successful SC or AMO still owns the bytes it wrote, so a later load does not
pair through it with an older store; a failed SC wrote nothing and changes
nothing. The stack pointer is read only for the stack-release diagnostic; it
never decides whether an access is a spill.

\section{Detector counters}

The detector writes these to \ct{riscv\_spill\_stats.txt}: one ROI block per
ROI and one global block.

\tablehead{3}{The counters the detector writes.}

{\small
\begin{xltabular}{\textwidth}{@{}C{40mm} L{1.0} L{1.0}@{}}
\toprule
\normalfont\textbf{Counter} & \textbf{Counted when} & \textbf{How to read it} \\
\midrule
\endfirsthead
\toprule
\normalfont\textbf{Counter} & \textbf{Counted when} & \textbf{How to read it} \\
\midrule
\endhead
\bottomrule
\endlastfoot
metric\_schema & Constant, 10 here.
  & Which spill definition produced the result. Results of schema 9 and
    earlier use other definitions and must not be mixed in. \\
total\_insts & Every committed instruction of the simulation.
  & Length of the whole program; the reference for \ct{roi/total}. \\
total\_load, total\_store & Every committed load or store.
  & Load and store counts of the whole program. \\
\midrule
roi\_total\_insts & Every instruction committed inside the ROI (the markers
  excluded).
  & Size of the measured region: the denominator of every ROI ratio. \\
roi\_load, roi\_store & Loads and stores committed inside the ROI.
  & All of them, spill or not. LR counts as a load and SC as a store; an AMO
    is neither. \\
roi\_spill\_store & A store meeting S1--S4, read back or not.
  & How often the program wrote a value to the stack: spills, callee saves,
    stack locals together. \\
roi\_spill\_event & A spill store read back at least once: one per store, at
  its first spill load.
  & How many of those stack writes were actually used again. Never above
    \ct{roi\_spill\_store} or \ct{roi\_spill\_load}. \\
roi\_spill\_load & A load meeting L1--L7. Repeats count.
  & How often the program read a stack-written value back. A slot read three
    times counts three, so it can exceed \ct{roi\_spill\_event}. \\
\midrule
roi\_\textit{p}\us l1\us hits, roi\_\textit{p}\us l1\us misses
  & L1D hits and misses of population \textit{p}'s accesses.
  & How much of that traffic L1D served, and how much went down to L2. \\
roi\_\textit{p}\us l2\us hits, roi\_\textit{p}\us l2\us misses
  & L2 hits and misses of population \textit{p}'s L1D misses.
  & The misses went on to DRAM: the most expensive accesses. \\
roi\_\textit{p}\us split\us insts
  & Instructions of \textit{p} whose access crosses a cache-line boundary;
    each makes two accesses.
  & Why L1 accesses can exceed the instruction count. Normally small. \\
roi\_spill\_event\us * & The store side of each event: the store's own
  outcome when it executed, added once at its first reload.
  & ``How did the stores that were read back fare in the cache?'' \\
\midrule
diag\us store\us map\us evicted
  & Store-table entries dropped at the size cap ($2^{20}$ entries).
  & Must be 0. Otherwise some pairs were lost and the spill counts are too
    low. \\
diag\us load\us after\us stack\us release
  & Spill loads of a byte that lay below \ct{sp} at some point after its
    store: the stack there had been released.
  & A suspicious pair: another frame may have reused the address. It does not
    change the main counters. \\
diag\us event\us after\us stack\us release
  & Spill events whose first spill load is such a load.
  & The same suspicion on the event side. \\
\end{xltabular}
}

\section{gem5 statistics the reports read}

Only the ROI block of \ct{stats.txt} is read. The statistics are dumped and
reset at the ROI boundaries.

\tablehead{4}{gem5 statistics used. L1D is \ct{system.cpu.dcache}, L2 is
\ct{system.l2}.}

{\small
\begin{xltabular}{\textwidth}{@{}C{52mm} L{0.75} L{1.25}@{}}
\toprule
\normalfont\textbf{Statistic} & \textbf{Used in} & \textbf{How to read it} \\
\midrule
\endfirsthead
\toprule
\normalfont\textbf{Statistic} & \textbf{Used in} & \textbf{How to read it} \\
\midrule
\endhead
\bottomrule
\endlastfoot
L1D demandHits, demandMisses, demandAccesses (\ct{::total})
  & \ct{L1D\_miss}
  & Overall success of L1D in this configuration: every load and store,
    spill or not. \\
L1D ReadReq.misses, ReadReq.accesses
  & \ct{other\_ld\_L1miss}
  & Ordinary loads. Removing the spill loads leaves the other loads. \\
L1D WriteReq.misses, WriteReq.accesses
  & \ct{other\_st\_L1miss}
  & Ordinary stores. Removing the spill stores leaves the other stores. \\
L2 demandHits, demandMisses
  & \ct{L2\_local\_miss}
  & Every request reaching L2: the L1D \emph{and} L1I misses. \\
L1D and L2 stackStoreHits, stackStoreMisses
  & Cross-check
  & Added by this project. The name is historical; the values must equal
    \ct{roi\_spill\_store\_l1/l2\_*}. \\
L1D and L2 spillLoadHits, spillLoadMisses
  & Cross-check
  & Added by this project; must equal \ct{roi\_spill\_load\_l1/l2\_*}. If they
    are equal, the detector and the caches saw the same accesses. \\
simInsts, numLoadInsts, numStoreInsts
  & Tests only
  & gem5's ROI block covers the detector's ROI: \ct{simInsts} $=$
    \ct{roi\_total\_insts} $+ 1$, the end marker. \\
\end{xltabular}
}

LR, SC and AMO accesses reach the cache as other request types, so they are in
neither \ct{ReadReq} nor \ct{WriteReq}, and so in no ``other'' rate.

\section{Values the reports compute}

\tablehead{5}{Derived values.}

{\small
\begin{tabularx}{\textwidth}{@{}C{44mm} L{0.9} L{1.1}@{}}
\toprule
\normalfont\textbf{Name} & \textbf{Formula} & \textbf{How to read it} \\
\midrule
roi\_spill\_insts & \ct{roi\_spill\_load} $+$ \ct{roi\_spill\_store}
  & All spill-like instructions. \\
roi\_total\_mem & \ct{roi\_load} $+$ \ct{roi\_store}
  & All memory instructions of the ROI. \\
roi\_unread\_spill\_store & \ct{roi\_spill\_store} $-$ \ct{roi\_spill\_event}
  & Stack writes never read back inside the ROI: overwritten first, dead, or
    still unread when the ROI ended. \\
roi\_non\_spill\_store & \ct{roi\_store} $-$ \ct{roi\_spill\_store}
  & Stores outside the stack (heap, globals), and SC. \\
\bottomrule
\end{tabularx}
}

\section{Program ratios}

They do not depend on the cache, so all five configurations show the same
values. Every ratio is a percentage.

\tablehead{6}{Program ratios.}

{\small
\begin{xltabular}{\textwidth}{@{}C{36mm} L{0.95} L{1.05}@{}}
\toprule
\normalfont\textbf{Ratio} & \textbf{Formula ($\times 100$)} & \textbf{How to read it} \\
\midrule
\endfirsthead
\toprule
\normalfont\textbf{Ratio} & \textbf{Formula ($\times 100$)} & \textbf{How to read it} \\
\midrule
\endhead
\bottomrule
\endlastfoot
roi/total & \ct{roi\_total\_insts} / \ct{total\_insts}
  & How much of the program was measured. Full runs: 95--100\%; low values
    mean the ROI may not represent the program. \\
load/roi & \ct{roi\_load} / \ct{roi\_total\_insts}
  & Read intensity of the program. \\
store/roi & \ct{roi\_store} / \ct{roi\_total\_insts}
  & Write intensity of the program. \\
spill/roi & (\ct{roi\_spill\_load} $+$ \ct{roi\_spill\_store}) /
  \ct{roi\_total\_insts}
  & \textbf{The main metric}: the share of instructions that write a value
    to the stack or read one back. High: register pressure, or many calls. \\
spill\_load/load & \ct{roi\_spill\_load} / \ct{roi\_load}
  & Share of the loads that read a stack-written value back. \\
spill\_store/store & \ct{roi\_spill\_store} / \ct{roi\_store}
  & Share of the stores that go to the stack, read back or not. \\
spill\_event/store & \ct{roi\_spill\_event} / \ct{roi\_store}
  & Share of the stores that write a value which is later read back. \\
spill\_event/ spill\_store & \ct{roi\_spill\_event} / \ct{roi\_spill\_store}
  & Reuse rate of stack writes. Near 100\%: almost every stack write is used
    again; low: many stack writes are never read. \\
spill/mem & (\ct{roi\_spill\_load} $+$ \ct{roi\_spill\_store}) /
  (\ct{roi\_load} $+$ \ct{roi\_store})
  & Share of the memory traffic that is spill-like. Unlike
    \ct{spill/roi}, the denominator holds memory instructions only. \\
\end{xltabular}
}

\section{Cache miss rates}

These depend on the configuration: the program and its instructions are the
same, only the hits and misses change. An \emph{access} is one request of a
load or store to a cache (a cross-line instruction makes two); a \emph{hit}
ends at that level; a \emph{miss} goes one level down (L1D $\to$ L2 $\to$
DRAM). The caches are write-allocate: a store miss first fetches the line.
An L2 rate is \emph{local}: its denominator is the L1D misses of that kind.

\tablehead{7}{Miss rates, in percent ($\times 100$). \textit{p} accesses
$=$ \ct{roi\_}\textit{p}\ct{\_l1\_hits} $+$
\ct{roi\_}\textit{p}\ct{\_l1\_misses}; likewise for L2.}

{\small
\begin{xltabular}{\textwidth}{@{}C{34mm} L{1.05} L{0.95}@{}}
\toprule
\normalfont\textbf{Rate} & \textbf{Formula} & \textbf{How to read it} \\
\midrule
\endfirsthead
\toprule
\normalfont\textbf{Rate} & \textbf{Formula} & \textbf{How to read it} \\
\midrule
\endhead
\bottomrule
\endlastfoot
L1D\_miss & L1D demand misses / L1D demand accesses
  & Overall L1D miss rate; the basic value to compare configurations. \\
spill\_ld\_L1miss & spill-load L1 misses / spill-load L1 accesses
  & Usually low: the slot was written shortly before, so its line is often
    still in L1D. \\
other\_ld\_L1miss & (ReadReq misses $-$ spill-load L1 misses) / (ReadReq
  accesses $-$ spill-load L1 accesses)
  & \textbf{Read next to} \ct{spill\_ld\_L1miss}: do spill loads miss less or
    more than ordinary loads? \\
spill\_st\_L1miss & spill-store L1 misses / spill-store L1 accesses
  & Miss rate of all stack writes. \\
spill\_event\_L1miss & spill-event L1 misses / spill-event L1 accesses
  & Miss rate of the stack writes that were read back. \\
unread\_st\_L1miss & (spill-store $-$ spill-event L1 misses) / (spill-store
  $-$ spill-event L1 accesses)
  & Miss rate of the stack writes never read back. \\
other\_st\_L1miss & (WriteReq misses $-$ spill-store L1 misses) / (WriteReq
  accesses $-$ spill-store L1 accesses)
  & Ordinary stores; compare with \ct{spill\_st\_L1miss}. \\
\midrule
L2\_local\_miss & L2 demand misses / L2 demand accesses
  & Share of the requests reaching L2 (L1D and L1I misses) that went on to
    DRAM. \\
spill\_ld\_L2local & spill-load L2 misses / spill-load L2 accesses
  & Of the spill loads that missed L1D, the share that missed L2 too. \\
spill\_st\_L2local & spill-store L2 misses / spill-store L2 accesses
  & The same for all stack writes. \\
spill\_event\_L2local & spill-event L2 misses / spill-event L2 accesses
  & The same for the stack writes read back. \\
unread\_st\_L2local & (spill-store $-$ spill-event L2 misses) / (spill-store
  $-$ spill-event L2 accesses)
  & The same for the stack writes never read back. \\
\end{xltabular}
}

The denominators of the L2 local rates can be tiny (a few accesses can give
66.7\%). Read them together with the absolute counts.

\section{Combining results}

\tablehead{8}{How several results become one number.}

{\small
\begin{tabularx}{\textwidth}{@{}l L{1.0} L{1.0}@{}}
\toprule
\textbf{Where} & \textbf{Formula} & \textbf{How to read it} \\
\midrule
SimPoint estimate
  & $C = N \cdot \sum_i w_i\, c_i / n_i \,\big/ \sum_i w_i$: $N$ the ROI
    length from the qemu profile, $n_i$ and $c_i$ the instructions and the
    counter of point $i$, $w_i$ its weight
  & A ref result is a weighted estimate, not a full run; its ratios are
    weighted per-instruction ratios. \\
pooled & $\sum$ numerators $/ \sum$ denominators over all completed runs
  & Large runs dominate: the share of the total traffic. \\
mean & Average of the per-benchmark values
  & Every benchmark weighs the same: the typical benchmark. \\
\bottomrule
\end{tabularx}
}

\section{Identities every result satisfies}

The test suite asserts these, and every completed result is checked against
them:
\begin{align*}
\ct{roi\_spill\_event} &\le \ct{roi\_spill\_store} \le \ct{roi\_store} \\
\ct{roi\_spill\_event} &\le \ct{roi\_spill\_load} \le \ct{roi\_load} \\
\ct{roi\_}p\ct{\_l1\_hits} + \ct{roi\_}p\ct{\_l1\_misses}
  &= \text{count of } p + \ct{roi\_}p\ct{\_split\_insts} \\
\ct{roi\_}p\ct{\_l2\_hits} + \ct{roi\_}p\ct{\_l2\_misses}
  &= \ct{roi\_}p\ct{\_l1\_misses} \\
\ct{roi\_spill\_event\_*} &\le \ct{roi\_spill\_store\_*} \\
\text{gem5 } \ct{stackStore*} = \ct{roi\_spill\_store\_*}, &\quad
\text{gem5 } \ct{spillLoad*} = \ct{roi\_spill\_load\_*}
\end{align*}

\section{Worked example}

\subsection*{A.\ Which instruction counts where}

The fragment runs inside the ROI; \ct{a2} and \ct{a3} point to the heap.

\tablehead{9}{Contribution of each instruction (empty: no change).}

{\small
\begin{tabularx}{\textwidth}{@{}>{\ttfamily}l c c c c c X@{}}
\toprule
\normalfont\textbf{Instruction} & \textbf{store} & \textbf{load}
  & \textbf{spill st.} & \textbf{event} & \textbf{spill ld.} & \textbf{Why} \\
\midrule
addi sp, sp, -32 & & & & & & Not a memory instruction \\
sd a0, 16(sp) & $+1$ & & $+1$ & & & Ordinary store to the stack \\
sd a1, 8(sp) & $+1$ & & $+1$ & & & Stack store; never read: unread \\
sw t0, 0(a2) & $+1$ & & & & & Heap store: a non-spill store \\
ld t1, 0(a3) & & $+1$ & & & & Heap load: an other load \\
ld a0, 16(sp) & & $+1$ & & $+1$ & $+1$ & Owner is the first \ct{sd}; its
  first reload: an event \\
ld a2, 16(sp) & & $+1$ & & & $+1$ & Same store read again: no new event \\
amoadd.w t2, t3, (a4) & & & & & & Atomic: neither load nor store \\
addi sp, sp, 32 & & & & & & Not a memory instruction \\
\midrule
\normalfont\textbf{Total} & 3 & 3 & 2 & 1 & 2 & \ct{roi\_total\_insts} $=$ 9 \\
\bottomrule
\end{tabularx}
}

Derived: \ct{roi\_unread\_spill\_store} $= 2 - 1 = 1$ (the slot
\ct{8(sp)}), \ct{roi\_non\_spill\_store} $= 3 - 2 = 1$ (the heap store).

\subsection*{B.\ From counters to ratios}

Assume a run reports \ct{total\_insts} $=$ 1000, \ct{total\_load} $=$ 300,
\ct{total\_store} $=$ 150; \ct{roi\_total\_insts} $=$ 800, \ct{roi\_load} $=$
240, \ct{roi\_store} $=$ 120; \ct{roi\_spill\_store} $=$ 60,
\ct{roi\_spill\_event} $=$ 45, \ct{roi\_spill\_load} $=$ 70.

\tablehead{10}{Program ratios of the example.}

{\small
\begin{tabularx}{\textwidth}{@{}>{\ttfamily}l l r X@{}}
\toprule
\normalfont\textbf{Ratio} & \textbf{Calculation} & \textbf{Value}
  & \textbf{Reading} \\
\midrule
roi/total & 800 / 1000 & 80.00\% & 80\% of the program was measured \\
spill/roi & (70 $+$ 60) / 800 & 16.25\% & About 16 of every 100
  instructions write to or read back from the stack \\
spill\_load/load & 70 / 240 & 29.17\% & Nearly a third of the loads are
  stack read-backs \\
spill\_store/store & 60 / 120 & 50.00\% & Half of the stores go to the
  stack \\
spill\_event/store & 45 / 120 & 37.50\% & 37.5\% of the stores write a value
  that is read later \\
spill\_event/spill\_store & 45 / 60 & 75.00\% & Three quarters of the stack
  writes are used again; 15 never are \\
spill/mem & 130 / 360 & 36.11\% & About a third of the memory traffic is
  spill-like \\
\bottomrule
\end{tabularx}
}

For the cache, assume no cross-line accesses and these outcomes:
spill stores 54 L1 hits and 6 misses, of which the spill events 42 and 3;
spill loads 68 and 2; L1D \ct{ReadReq} 240 accesses and 30 misses,
\ct{WriteReq} 120 and 18; at L2, spill stores 4 hits and 2 misses (events 3
and 0), spill loads 1 and 1; L2 in total 60 accesses (48 from L1D, 12 from
L1I) and 15 misses.

\tablehead{11}{Miss rates of the example.}

{\small
\begin{tabularx}{\textwidth}{@{}>{\ttfamily}l l r X@{}}
\toprule
\normalfont\textbf{Rate} & \textbf{Calculation} & \textbf{Value}
  & \textbf{Reading} \\
\midrule
spill\_ld\_L1miss & 2 / (68 $+$ 2) & 2.86\% & Spill loads almost always hit
  \ldots \\
other\_ld\_L1miss & (30 $-$ 2) / (240 $-$ 70) & 16.47\% & \ldots\ ordinary
  loads miss far more: spill traffic is cache friendly \\
spill\_st\_L1miss & 6 / (54 $+$ 6) & 10.00\% & All stack writes \\
spill\_event\_L1miss & 3 / (42 $+$ 3) & 6.67\% & Stack writes read back \\
unread\_st\_L1miss & (6 $-$ 3) / (60 $-$ 45) & 20.00\% & Stack writes never
  read miss more \\
other\_st\_L1miss & (18 $-$ 6) / (120 $-$ 60) & 20.00\% & Ordinary stores \\
L1D\_miss & 48 / 360 & 13.33\% & All L1D accesses \\
L2\_local\_miss & 15 / 60 & 25.00\% & All L2 accesses \\
spill\_ld\_L2local & 1 / (1 $+$ 1) & 50.00\% & Denominator 2: too few to
  interpret \\
unread\_st\_L2local & (2 $-$ 0) / (6 $-$ 3) & 66.67\% & Denominator 3: too
  few to interpret \\
\bottomrule
\end{tabularx}
}

{\raggedright Checks: $45 \le 60 \le 120$, $45 \le 70 \le 240$, $54 + 6 = 60$ and
$4 + 2 = 6$ hold. At L1D gem5 must report \ct{stackStoreHits} $=$ 54,
\ct{stackStoreMisses} $=$ 6, \ct{spillLoadHits} $=$ 68,
\ct{spillLoadMisses} $=$ 2, and at L2 4, 2, 1 and 1.\par}

\section{Implementation and verification}

\begin{tabularx}{\textwidth}{@{}X l@{}}
\toprule
\textbf{Component} & \textbf{Source} \\
\midrule
Rules and counters
  & \texttt{src/cpu/simple/spill\_detector.\{hh,cc\}} \\
Commit hook, request tags, LR/SC/AMO handling
  & \texttt{src/cpu/simple/timing.cc} \\
SimPoint ROI (the interval)
  & \texttt{configs/common/Simulation.py} \\
Per-population cache counters
  & \texttt{src/mem/cache/base.\{hh,cc\}} \\
Ratios and miss rates
  & \texttt{cache\_scripts/cache\_status.py} \\
SimPoint estimates
  & \texttt{cache\_scripts/simpoint\_combine.py} \\
\bottomrule
\end{tabularx}

The rules are covered by 27 hand-verified RISC-V assembly test cases in
\texttt{tests/spill\_test/}: partial overwrites, atomics, a SimPoint interval,
frame reuse and stack growth among them. They assert every identity above and
that the gem5 cache statistics and the detector counters agree at both cache
levels. On 12 November 2026 every completed schema-10 result of the campaign,
170 test runs and 3405 SimPoint points, satisfied every identity;
\ct{diag\_store\_map\_evicted} was 0 everywhere, and
\ct{diag\_load\_after\_stack\_release} was 0.05\% of the spill loads.

% ---- Turkish version of the worked example (section 10) ----
% The engine's T1 fonts lack the Turkish letters as Unicode input, so the
% text uses LaTeX accent commands (\c{s}, \u{g}, {\i}, \.{I}, ...). English
% hyphenation patterns would split Turkish words wrongly: hyphenation is off.
\clearpage
\begingroup
\hyphenpenalty=10000 \exhyphenpenalty=10000 \emergencystretch=3em

\section{\"Ornek (T\"urk\c{c}e)}

Bu b\"ol\"um, 10.\ b\"ol\"umdeki \"orne\u{g}in T\"urk\c{c}e anlat{\i}m{\i}d{\i}r.

\subsection*{A.\ Hangi komut hangi sayac{\i} art{\i}r{\i}r?}

Kod par\c{c}as{\i} ROI i\c{c}inde \c{c}al{\i}\c{s}{\i}yor; \ct{a2} ve \ct{a3} heap'i g\"osteriyor.

\tablehead{12}{Her komutun saya\c{c}lara katk{\i}s{\i} (bo\c{s}: de\u{g}i\c{s}mez).}

{\small
\begin{tabularx}{\textwidth}{@{}>{\ttfamily}l c c c c c X@{}}
\toprule
\normalfont\textbf{Komut} & \textbf{store} & \textbf{load}
  & \textbf{spill st.} & \textbf{event} & \textbf{spill ld.} & \textbf{Neden} \\
\midrule
addi sp, sp, -32 & & & & & & Bellek komutu de\u{g}il \\
sd a0, 16(sp) & $+1$ & & $+1$ & & & Stack'e normal store \\
sd a1, 8(sp) & $+1$ & & $+1$ & & & Stack'e store; hi\c{c} okunmayacak: unread \\
sw t0, 0(a2) & $+1$ & & & & & Heap'e yaz{\i}yor: non-spill store \\
ld t1, 0(a3) & & $+1$ & & & & Heap'ten okuyor: other load \\
ld a0, 16(sp) & & $+1$ & & $+1$ & $+1$ & Byte'{\i}n sahibi ilk \ct{sd}; ilk
  reload oldu\u{g}u i\c{c}in event \\
ld a2, 16(sp) & & $+1$ & & & $+1$ & Ayn{\i} store'un ikinci okumas{\i}: yeni
  event yok \\
amoadd.w t2, t3, (a4) & & & & & & Atomic: hi\c{c}bir load/store sayac{\i}na
  girmez \\
addi sp, sp, 32 & & & & & & Bellek komutu de\u{g}il \\
\midrule
\normalfont\textbf{Toplam} & 3 & 3 & 2 & 1 & 2 & \ct{roi\_total\_insts} $=$ 9 \\
\bottomrule
\end{tabularx}
}

T\"uretilenler: \ct{roi\_unread\_spill\_store} $= 2 - 1 = 1$ (\ct{8(sp)}
slot'u), \ct{roi\_non\_spill\_store} $= 3 - 2 = 1$ (heap store'u).

\subsection*{B.\ Say{\i}larla oranlar}

Bir ko\c{s}unun \c{s}u sonu\c{c}lar{\i} verdi\u{g}ini d\"u\c{s}\"unelim:
\begin{itemize}
\setlength{\itemsep}{1pt}\setlength{\topsep}{2pt}
\item T\"um program: \ct{total\_insts} $=$ 1000, \ct{total\_load} $=$ 300,
  \ct{total\_store} $=$ 150
\item ROI: \ct{roi\_total\_insts} $=$ 800, \ct{roi\_load} $=$ 240,
  \ct{roi\_store} $=$ 120
\item Spill saya\c{c}lar{\i}: \ct{roi\_spill\_store} $=$ 60,
  \ct{roi\_spill\_event} $=$ 45, \ct{roi\_spill\_load} $=$ 70
\end{itemize}

\tablehead{13}{\"Orne\u{g}in program oranlar{\i}.}

{\small
\begin{tabularx}{\textwidth}{@{}>{\ttfamily}l l r X@{}}
\toprule
\normalfont\textbf{Oran} & \textbf{Hesap} & \textbf{Sonu\c{c}} & \textbf{Yorum} \\
\midrule
roi/total & 800 / 1000 & \%80,00 & Program{\i}n \%80'ini \"ol\c{c}t\"uk \\
spill/roi & (70 $+$ 60) / 800 & \textbf{\%16,25} & Her 100 komuttan yakla\c{s}{\i}k
  16's{\i} stack'e yazma ya da stack'ten geri okuma \\
spill\_load/load & 70 / 240 & \%29,17 & Load'lar{\i}n yakla\c{s}{\i}k \"u\c{c}te biri
  stack'ten geri okuma \\
spill\_store/store & 60 / 120 & \%50,00 & Store'lar{\i}n yar{\i}s{\i} stack'e gidiyor \\
spill\_event/store & 45 / 120 & \%37,50 & Store'lar{\i}n \%37,5'i stack'e
  yaz{\i}l{\i}p sonra okunan de\u{g}er \\
spill\_event/spill\_store & 45 / 60 & \%75,00 & Stack'e yaz{\i}lanlar{\i}n d\"ortte
  \"u\c{c}\"u geri kullan{\i}lm{\i}\c{s}; 15'i (unread) hi\c{c} okunmam{\i}\c{s} \\
spill/mem & 130 / 360 & \%36,11 & Bellek trafi\u{g}inin yakla\c{s}{\i}k \"u\c{c}te biri
  spill-like \\
\bottomrule
\end{tabularx}
}

Cache taraf{\i} i\c{c}in cache line s{\i}n{\i}r{\i}n{\i} a\c{s}an access olmad{\i}\u{g}{\i}n{\i}
(\ct{split\_insts} $=$ 0) ve \c{s}u sonu\c{c}lar{\i} varsayal{\i}m:
\begin{itemize}
\setlength{\itemsep}{1pt}\setlength{\topsep}{2pt}
\item Spill store'lar: L1'de 54 hit, 6 miss; bunlardan spill event olanlar
  42 hit, 3 miss
\item Spill load'lar: L1'de 68 hit, 2 miss
\item L1D \ct{ReadReq}: 240 access, 30 miss; \ct{WriteReq}: 120 access, 18
  miss
\item L2: spill store'lar 4 hit, 2 miss (event'ler 3 hit, 0 miss); spill
  load'lar 1 hit, 1 miss
\item L2 toplam{\i}: 60 access (48'i L1D'den, 12'si L1I'dan), 15 miss
\end{itemize}

\tablehead{14}{\"Orne\u{g}in miss oranlar{\i}.}

{\small
\begin{tabularx}{\textwidth}{@{}>{\ttfamily}l l r X@{}}
\toprule
\normalfont\textbf{Oran} & \textbf{Hesap} & \textbf{Sonu\c{c}} & \textbf{Yorum} \\
\midrule
spill\_ld\_L1miss & 2 / (68 $+$ 2) & \textbf{\%2,86} & Spill load'lar
  neredeyse hep hit ediyor \ldots \\
other\_ld\_L1miss & (30 $-$ 2) / (240 $-$ 70) & \textbf{\%16,47} & \ldots\
  normal load'lar \c{c}ok daha fazla miss ediyor \\
spill\_st\_L1miss & 6 / (54 $+$ 6) & \%10,00 & B\"ut\"un stack yazmalar{\i} \\
spill\_event\_L1miss & 3 / (42 $+$ 3) & \%6,67 & Geri okunan stack
  yazmalar{\i} \\
unread\_st\_L1miss & (6 $-$ 3) / (60 $-$ 45) & \%20,00 & Hi\c{c} okunmayanlar
  daha \c{c}ok miss ediyor \\
other\_st\_L1miss & (18 $-$ 6) / (120 $-$ 60) & \%20,00 & Spill olmayan
  store'lar \\
L1D\_miss & 48 / 360 & \%13,33 & B\"ut\"un L1D access'leri \\
L2\_local\_miss & 15 / 60 & \%25,00 & B\"ut\"un L2 access'leri \\
spill\_ld\_L2local & 1 / (1 $+$ 1) & \%50,00 & Payda 2: yorumlamak i\c{c}in \c{c}ok
  az \\
unread\_st\_L2local & (2 $-$ 0) / (6 $-$ 3) & \%66,67 & Payda 3: yorumlamak
  i\c{c}in \c{c}ok az \\
\bottomrule
\end{tabularx}
}

\textbf{Okumas{\i}:}
\begin{itemize}
\setlength{\itemsep}{1pt}\setlength{\topsep}{2pt}
\item Spill load'lar L1D'de \c{c}ok az miss ediyor: di\u{g}er load'lardaki
  \%16,47'ye kar\c{s}{\i} \%2,86. Yani spill trafi\u{g}i cache dostu; stack'e yaz{\i}lan
  de\u{g}er geri okunurken genelde h\^al\^a L1'de.
\item Hi\c{c} okunmayan spill store'lar (\%20) okunanlardan (\%6,67) daha \c{c}ok
  miss ediyor. L2 local oranlar{\i} ise bu kadar k\"u\c{c}\"uk paydalarla
  yorumlanmamal{\i}.
\item Kontrol: $45 \le 60 \le 120$, $45 \le 70 \le 240$, $54 + 6 = 60$ ve
  $4 + 2 = 6$ tutuyor. gem5'in L1D'de \ct{stackStoreHits} $=$ 54,
  \ct{stackStoreMisses} $=$ 6, \ct{spillLoadHits} $=$ 68,
  \ct{spillLoadMisses} $=$ 2; L2'de 4, 2, 1 ve 1 g\"ostermesi gerekir.
\end{itemize}

\subsection*{C.\ Ger\c{c}ek veriden iki u\c{c}}

Test girdisi, small-2way konfig\"urasyonu:
\begin{itemize}
\setlength{\itemsep}{1pt}\setlength{\topsep}{2pt}
\item \textbf{cactuBSSN:} \ct{spill/roi} \%33,48,
  \ct{spill\_event/spill\_store} \%99,36. Komutlar{\i}n \"u\c{c}te biri spill-like
  ve stack'e yaz{\i}lan neredeyse her \c{s}ey geri okunuyor; register bask{\i}s{\i} \c{c}ok
  y\"uksek.
\item \textbf{parest:} \ct{spill/roi} \%2,13. Spill trafi\u{g}i az.
\end{itemize}

\endgroup

\end{document}
